\subsection{DP en Árboles con Recorrido Iterativo (emparejamiento máximo)}
\label{alg:6-6}

\graybold{Preguntas guía:}

\begin{itemize}
\item ¿El grafo es un ÁRBOL y me piden optimizar algo que se puede decidir subárbol por subárbol?
\item ¿El estado de cada nodo se resume en pocos casos (usado / libre, tomado / no tomado, pintado de tal color)?
\item ¿La respuesta de un nodo depende solo de las respuestas de sus hijos, permitiendo procesar de las hojas hacia la raíz?
\item ¿La profundidad puede llegar a \ten{5} o más, obligando a un recorrido iterativo en vez de recursivo?
\end{itemize}

\graybold{Utilidad:} Es la plantilla del DP sobre árboles: se recorre el árbol de las hojas hacia la raíz y cada nodo combina los resultados de sus hijos según un estado pequeño. Con dos estados por nodo se resuelven emparejamiento máximo, conjunto independiente máximo (con o sin pesos), cobertura mínima de vértices y coloreos; con más estados, variantes como "a lo sumo $k$ hijos elegidos". Todas comparten la estructura de esta entrada: un orden que garantiza hijos antes que padres, y una combinación aditiva con una corrección por el caso especial.

\graybold{Complejidad:} \bigO{n} en tiempo y memoria: cada nodo y cada arista se visitan una vez.

\graybold{Fundamento teórico:} Para cada nodo $v$ se definen dos valores sobre su subárbol: $dp_0[v]$, el mejor emparejamiento dejando a $v$ SIN emparejar, y $dp_1[v]$, el mejor dejándolo emparejado con alguno de sus hijos. Si $v$ queda libre, cada hijo aporta lo mejor que pueda por su cuenta: $dp_0[v] = \sum_{c} \max(dp_0[c], dp_1[c])$. Si $v$ se empareja con un hijo $c$, ese hijo debe estar libre, así que aporta $dp_0[c] + 1$ en lugar de $\max(dp_0[c], dp_1[c])$, y el resto de los hijos sigue aportando su máximo; por lo tanto
\[dp_1[v] = dp_0[v] + \max_{c}\big(dp_0[c] + 1 - \max(dp_0[c], dp_1[c])\big).\]
Escrito así, no hace falta recorrer los hijos dos veces ni guardar listas de hijos: basta acumular en el padre, mientras se procesan los hijos, la suma de los máximos y por separado la mayor de esas diferencias. La respuesta del nodo es $\max(dp_0[v], dp_1[v])$, que equivale a sumar la diferencia solo si es positiva. Hay dos detalles de implementación que evitan problemas. Primero, el ORDEN: se construye un recorrido por capas desde la raíz, que garantiza que todo nodo aparece después de su padre; recorrerlo AL REVÉS garantiza que los hijos se procesan antes, que es lo que el DP necesita, y evita la recursión, inviable con profundidad $2\cdot10^5$. Segundo, no hace falta separar hijos de padre en la lista de adyacencia: al marcar el padre de cada nodo durante la construcción del orden, el recorrido inverso ya sabe a quién aportar.

\graybold{Ejercicio aplicado}

(CSES 1130 — Tree Matching) Dado un árbol de $n$ nodos, hallar la máxima cantidad de aristas que se pueden elegir sin que dos compartan un nodo (emparejamiento máximo).

\graybold{I/O:} $n \le 2\cdot10^5$ con dos enteros por arista $\Rightarrow$ lectura total + tokenización (patrón 2), separando extremos con slices de paso 2. Se atiende aparte $n = 1$, donde no hay aristas que leer. Verificado contra el caso de ejemplo y contrastado con una fuerza bruta que prueba todos los subconjuntos de aristas en 500 árboles de hasta 11 nodos (caminos, estrellas, orugas, binarios y aleatorios, con los nodos renumerados al azar y las aristas mezcladas), sin discrepancias. Con $n = 2\cdot10^5$ tarda entre \aprox{}0.19s y \aprox{}0.26s en todas las formas medidas, incluido un camino de profundidad $2\cdot10^5$ donde una versión recursiva desbordaría la pila; el límite es de 1 segundo.

\graybold{Código solución}

\begin{lstlisting}[language=Python]
import sys

def solve():
    data = sys.stdin.buffer.read().split()
    n = int(data[0])
    if n == 1:
        print(0); return
    adj = [[] for _ in range(n + 1)]
    fin = 1 + 2*(n - 1)
    for a, b in zip(map(int, data[1:fin:2]), map(int, data[2:fin:2])):
        adj[a].append(b)
        adj[b].append(a)
    # orden por capas: cada nodo aparece despues de su padre
    padre = [0]*(n + 1)
    orden = [1]; padre[1] = 1          # la raiz se marca a si misma para no revisitarla
    for u in orden:                    # la lista crece mientras se recorre
        for v in adj[u]:
            if not padre[v]:
                padre[v] = u
                orden.append(v)
    libre = [0]*(n + 1)   # dp0: mejor emparejamiento del subarbol con el nodo LIBRE
    ganan = [0]*(n + 1)   # mayor (dp0[hijo] + 1 - max(dp0,dp1)[hijo]) entre los hijos
    for v in reversed(orden):          # al revés: los hijos ya estan procesados
        # max(dp0[v], dp1[v]): emparejar a v con un hijo solo si la diferencia suma
        mejor = libre[v] + (ganan[v] if ganan[v] > 0 else 0)
        if v == 1:
            print(mejor); return
        p = padre[v]
        libre[p] += mejor              # aporte de v si el padre queda libre
        # si el padre se empareja CON v, v debe quedar libre: aporta dp0[v] + 1
        cand = libre[v] + 1 - mejor
        if cand > ganan[p]:
            ganan[p] = cand

solve()
\end{lstlisting}
